% =============================================================================
%  Football ROI Pipeline — Rapport final
%  Nassim El Kharroubi — ISEP, École d'ingénieurs du numérique
%  Compilation : pdflatex RAPPORT.tex (2 passes pour le sommaire)
%                OU  xelatex RAPPORT.tex (recommandé)
%
%  Structure des fichiers attendue :
%    docs/RAPPORT.tex          (ce fichier)
%    screens/logo.jpg
%    screens/airflow.png
%    screens/arborescence_datalake.png
%    screens/top_clubs_roi.png
%    screens/roi_par_ligue.png
%    screens/evolution_roi_barca_psg.png
%    screens/budget_vs_points.png
%    screens/budget_vs_stage_cl.png
%    screens/evolution_cl_grands_clubs.png
%    screens/data_table_roi.png
%    screens/roi_par_cluster.png
% =============================================================================

\documentclass[11pt,a4paper]{article}

% --- Encodage et langue ---
\usepackage[utf8]{inputenc}
\usepackage[T1]{fontenc}
\usepackage[french]{babel}
\usepackage{lmodern}

% --- Mise en page ---
\usepackage[margin=2cm,top=2cm,bottom=2.5cm]{geometry}
\usepackage{setspace}
\setstretch{1.15}
\setlength{\parindent}{0pt}
\setlength{\parskip}{0.4em}

% --- Graphiques et tableaux ---
\usepackage{graphicx}
\graphicspath{{../screens/}{./screens/}{../../screens/}}
\usepackage{booktabs}
\usepackage{array}
\usepackage{tabularx}
\usepackage{caption}
\captionsetup{font=small,labelfont=bf,format=hang}

% --- Couleurs ---
\usepackage{xcolor}
\definecolor{isepblue}{RGB}{0,82,156}
\definecolor{lightgray}{RGB}{245,245,245}

% --- Titres de sections en couleur ---
\usepackage{titlesec}
\titleformat{\section}
  {\Large\bfseries\color{isepblue}}{\thesection.}{0.6em}{}
\titleformat{\subsection}
  {\large\bfseries\color{isepblue!85}}{\thesubsection}{0.6em}{}

% --- En-tête et pied de page ---
\usepackage{fancyhdr}
\pagestyle{fancy}
\fancyhf{}
\fancyhead[L]{\small\textit{Football ROI Pipeline}}
\fancyhead[R]{\small Nassim El Kharroubi}
\fancyfoot[C]{\thepage}
\renewcommand{\headrulewidth}{0.4pt}

% --- Liens hypertexte ---
\usepackage{hyperref}
\hypersetup{
  colorlinks=true,
  linkcolor=isepblue,
  urlcolor=isepblue,
  citecolor=isepblue,
  pdftitle={Football ROI Pipeline},
  pdfauthor={Nassim El Kharroubi}
}

% --- Listings (code) ---
\usepackage{listings}
\lstset{
  basicstyle=\ttfamily\small,
  backgroundcolor=\color{lightgray},
  frame=single,
  framesep=4pt,
  rulecolor=\color{gray!30},
  breaklines=true,
  showstringspaces=false,
  columns=flexible,
  language=Python
}

% =============================================================================
\begin{document}

% --- Page de garde ---------------------------------------------------------
\thispagestyle{empty}
\begin{center}
  \vspace*{1.5cm}
  \includegraphics[width=0.45\textwidth]{logo.jpg}\\[2.5cm]

  {\Huge\bfseries\color{isepblue} Football ROI Pipeline}\\[0.6cm]
  {\Large Analyse du retour sur investissement sportif\\des clubs européens}\\[2cm]

  \rule{0.7\textwidth}{0.4pt}\\[0.6cm]
  {\large \textbf{Projet Apache Airflow \& Big Data}}\\[0.4cm]
  4\textsuperscript{ème} année — École d'ingénieurs du numérique\\[2cm]

  \begin{tabular}{rl}
    \textbf{Auteur} & : Nassim El Kharroubi \\
    \textbf{Date}   & : Juin 2026 \\
    \textbf{Code source} & : \href{https://github.com/Nassim821/football-roi-pipeline}{github.com/Nassim821/football-roi-pipeline}
  \end{tabular}
  \vfill
\end{center}
\newpage

% --- Sommaire --------------------------------------------------------------
\thispagestyle{empty}
\tableofcontents
\newpage
\setcounter{page}{1}

% =============================================================================
\section{Introduction et problématique}

\subsection*{Contexte}

Le football européen vit une période d'inflation historique. Les clubs dépensent
chaque année des sommes record en transferts et en salaires, mais cet argent ne
se traduit pas toujours en performances sportives. Certains clubs au budget
modeste réussissent à finir devant des géants financiers, et d'autres dépensent
des centaines de millions d'euros sans même gagner un trophée.

\subsection*{Question business}

Le projet répond à une question simple mais rarement posée dans les classements
traditionnels :

\begin{center}
\textit{« Quels clubs investissent le mieux leur argent ? »}
\end{center}

Pour y répondre, j'ai inventé une métrique métier : le \textbf{ROI sportif},
défini comme le nombre de points obtenus en championnat divisé par la valeur de
l'effectif en millions d'euros.

\[
\text{ROI} = \frac{\text{points du championnat}}{\text{valeur de l'effectif (M\texteuro)}}
\]

Plus un club obtient de points avec un petit effectif, plus son ROI est élevé :
il est « efficace ». À l'inverse, un club avec un énorme effectif mais peu de
points a un ROI faible : il « gaspille ».

\subsection*{Périmètre}

\begin{itemize}
  \item \textbf{5 championnats majeurs} : Ligue 1, Premier League, La Liga, Bundesliga, Serie A
  \item \textbf{La Champions League} (avec un ROI adapté basé sur le stade atteint)
  \item \textbf{3 saisons} : 2022, 2023 et 2024
\end{itemize}

\subsection*{Apport du machine learning}

En plus du ROI, j'utilise un algorithme de clustering (K-Means avec 4 groupes)
pour identifier automatiquement des profils de clubs : les \textit{Efficaces},
les \textit{Dominants}, les \textit{Standards} et les \textit{Gaspilleurs}.

% =============================================================================
\section{Architecture technique}

\subsection*{Vue d'ensemble}

Le projet suit le pattern \textbf{Medallion}, très courant en data engineering.
Il consiste à organiser le data lake en trois couches successives : \textit{raw}
(données brutes), \textit{formatted} (données nettoyées) et \textit{usage}
(données prêtes à l'emploi).

\subsection*{Stack utilisée}

\begin{table}[h]
\centering
\small
\begin{tabularx}{\textwidth}{@{}lllX@{}}
\toprule
\textbf{Étape} & \textbf{Technologie} & \textbf{Rôle} \\
\midrule
Ingestion       & Python (requests, cloudscraper) & APIs REST + scraping HTML \\
Stockage        & JSON puis Parquet               & Brut traçable / typé et compact \\
Compute         & PySpark + pandas                & Jointures Spark, formatage pandas \\
Orchestration   & Apache Airflow 3.2.1            & Dépendances, retries, planification \\
ML              & scikit-learn (K-Means)          & Clustering non supervisé \\
Indexation      & Elasticsearch 9.0.0             & Stockage analytique et recherche \\
Visualisation   & Kibana 9.0.0                    & Dashboard interactif \\
Conteneurisation & Docker Compose                 & Reproductibilité de l'environnement \\
\bottomrule
\end{tabularx}
\caption{Récapitulatif des technologies utilisées.}
\end{table}

\subsection*{Pourquoi ces choix ?}

\begin{itemize}
  \item \textbf{Airflow} plutôt qu'un cron : pour gérer automatiquement les
        dépendances entre tâches, relancer en cas d'erreur, et visualiser
        l'avancement de chaque exécution.
  \item \textbf{Parquet} plutôt que CSV ou JSON : format colonnaire compressé,
        environ 10× plus rapide à lire et beaucoup plus léger.
  \item \textbf{Spark} plutôt que pandas pour les jointures : Spark gère mieux
        les volumes importants tout en restant lisible.
  \item \textbf{Elasticsearch + Kibana} plutôt que Tableau ou PowerBI : open
        source, performant pour la recherche, et intégré nativement.
\end{itemize}

% =============================================================================
\section{Sources de données et data lake}

\subsection*{Trois sources combinées}

Pour calculer le ROI, j'ai eu besoin de deux types d'informations : les
performances sportives (points, rang) et les valeurs marchandes (budget de
l'effectif).

\begin{table}[h]
\centering
\small
\begin{tabularx}{\textwidth}{@{}llXl@{}}
\toprule
\textbf{Source} & \textbf{Type} & \textbf{Donnée récupérée} & \textbf{Volume / run} \\
\midrule
API-Football (standings) & API REST officielle & Classements et résultats des 5 ligues & 15 appels \\
Transfermarkt            & Web scraping HTML   & Valeurs marchandes des effectifs       & 15 scrapes \\
API-Football (CL)        & API REST officielle & Matchs Champions League                & 3 appels \\
\bottomrule
\end{tabularx}
\caption{Sources de données ingérées par le pipeline.}
\end{table}

\subsection*{Organisation du data lake}

Le data lake respecte une convention de nommage stricte :
\texttt{data/<layer>/<group>/<dataEntity>/<dateVersion>/}.

\begin{figure}[h]
\centering
\includegraphics[width=0.55\textwidth]{arborescence_datalake.png}
\caption{Arborescence du data lake en 3 couches : raw, formatted et usage.}
\end{figure}

\subsection*{Volumétrie finale}

Après le pipeline complet, j'obtiens \textbf{312 lignes business} indexées dans
Elasticsearch :

\begin{itemize}
  \item 254 lignes pour les championnats domestiques (5 ligues × $\sim$17 clubs × 3 saisons)
  \item 58 lignes pour la Champions League
  \item 254 lignes enrichies par le clustering ML
\end{itemize}

% =============================================================================
\section{Le pipeline Airflow}

\subsection*{Vue d'ensemble du DAG}

Le DAG \texttt{football\_roi\_pipeline} est composé de \textbf{10 tâches}
organisées en 5 étapes successives. Il tourne automatiquement tous les jours à
6h UTC.

\begin{figure}[h]
\centering
\includegraphics[width=\textwidth]{airflow.png}
\caption{Vue Graph du DAG Airflow : 10 tâches organisées en 5 étapes successives.}
\end{figure}

\subsection*{Détail des 10 tâches}

\begin{table}[h]
\centering
\small
\begin{tabularx}{\textwidth}{@{}llX@{}}
\toprule
\textbf{Tâche} & \textbf{Couche} & \textbf{Rôle} \\
\midrule
\texttt{source\_to\_raw\_1} & source $\to$ raw       & Récupère les classements via l'API-Football \\
\texttt{source\_to\_raw\_2} & source $\to$ raw       & Scrape Transfermarkt pour les valeurs marchandes \\
\texttt{source\_to\_raw\_3} & source $\to$ raw       & Récupère les matchs de Champions League \\
\texttt{raw\_to\_formatted\_1} & raw $\to$ formatted & Parse les JSON football $\to$ Parquet typé \\
\texttt{raw\_to\_formatted\_2} & raw $\to$ formatted & Nettoie Transfermarkt (ex : "\texteuro 45.20m" $\to$ 45.2) \\
\texttt{raw\_to\_formatted\_3} & raw $\to$ formatted & Calcule le stade max atteint en CL par club \\
\texttt{produce\_usage}        & formatted $\to$ usage & JOIN Spark + calcul du ROI sportif \\
\texttt{produce\_usage\_cl}    & formatted $\to$ usage & JOIN CL + calcul du \texttt{cl\_roi} \\
\texttt{ml\_clustering}        & usage $\to$ usage     & K-Means 4 clusters et labelling automatique \\
\texttt{index\_to\_elastic}    & usage $\to$ ES        & Pousse les 3 datasets vers Elasticsearch \\
\bottomrule
\end{tabularx}
\caption{Liste des 10 tâches du DAG et leur rôle.}
\end{table}

\subsection*{Pourquoi 3 couches ?}

Cette structure est une \textit{best practice} de data engineering :

\begin{itemize}
  \item La couche \textbf{raw} est immuable. Si demain l'API change ou tombe,
        j'ai toujours mes données originales pour recalculer.
  \item La couche \textbf{formatted} est techniquement propre (types corrects,
        dates UTC, doublons supprimés) mais ne contient pas encore de logique métier.
  \item La couche \textbf{usage} est spécifique à mon cas d'usage (le calcul
        du ROI). Plusieurs cas d'usage peuvent cohabiter à ce niveau.
\end{itemize}

% =============================================================================
\section{Le machine learning}

\subsection*{Choix de l'algorithme}

J'ai choisi \textbf{K-Means} parce que :

\begin{itemize}
  \item Le problème est non supervisé : je n'ai pas de labels au départ, je
        veux découvrir des groupes naturels dans les données.
  \item C'est un algorithme simple à interpréter : chaque cluster est défini
        par un centroïde, facile à expliquer.
  \item Il est rapide et léger : 254 lignes traitées en moins d'une seconde.
\end{itemize}

\subsection*{Features et préparation}

Trois variables servent au clustering :
\texttt{squad\_value\_m} (budget), \texttt{roi} (efficacité) et \texttt{points}
(résultats). Elles sont normalisées avec \texttt{StandardScaler} pour éviter
qu'une variable à grande échelle ne domine l'algorithme.

\subsection*{Choix de $k=4$}

J'ai fixé le nombre de clusters à 4 pour obtenir 4 archétypes métiers
correspondant aux quadrants du plan (budget, ROI) :

\begin{itemize}
  \item Petit budget + gros ROI $\to$ \textbf{Efficaces}
  \item Gros budget + gros ROI $\to$ \textbf{Dominants}
  \item Gros budget + petit ROI $\to$ \textbf{Gaspilleurs}
  \item Petit budget + petit ROI $\to$ \textbf{Standards}
\end{itemize}

\subsection*{Logique de labelling}

Plutôt que d'attribuer les labels à la main, j'ai écrit une fonction qui
calcule la distance euclidienne entre chaque centroïde et les 4 archétypes
dans le plan normalisé. Une \textbf{assignation 1-pour-1} garantit que chaque
label est utilisé exactement une fois (la première version avait un bug à ce
sujet, expliqué en section 7).

\subsection*{Résultats du clustering}

\begin{table}[h]
\centering
\small
\begin{tabular}{lcccc}
\toprule
\textbf{Cluster} & \textbf{Budget moyen (M\texteuro)} & \textbf{ROI moyen} & \textbf{Points moyens} & \textbf{Nombre de clubs} \\
\midrule
Efficaces    & 85  & 0.53 & 43 & 119 \\
Dominants    & 155 & 0.27 & 40 & 31 \\
Standards    & 423 & 0.15 & 59 & 26 \\
Gaspilleurs  & 980 & 0.08 & 79 & 104 \\
\bottomrule
\end{tabular}
\caption{Caractéristiques des 4 clusters identifiés par K-Means.}
\end{table}

% =============================================================================
\section{Résultats : analyse des 8 visualisations}

Le dashboard Kibana regroupe 8 visualisations. Chacune répond à une question
spécifique du projet.

% --- Viz 1 ---
\subsection*{Viz 1 — Top des clubs par ROI}

\begin{figure}[h]
\centering
\includegraphics[width=0.85\textwidth]{top_clubs_roi.png}
\caption{Top 15 des clubs par ROI moyen sur les 3 saisons.}
\end{figure}

\textbf{Lecture :} Clermont Foot domine largement avec un ROI proche de 0.82,
suivi de Rayo Vallecano, Holstein Kiel et Cadiz.

\textbf{Insight :} Ce sont tous des clubs à petit budget. Le ROI inverse
complètement le palmarès classique : les « petits » deviennent les plus
efficaces, parce qu'ils marquent beaucoup de points par rapport à leur faible
investissement.

% --- Viz 2 ---
\subsection*{Viz 2 — ROI moyen par ligue}

\begin{figure}[h]
\centering
\includegraphics[width=0.85\textwidth]{roi_par_ligue.png}
\caption{ROI moyen par championnat sur les 3 saisons.}
\end{figure}

\textbf{Lecture :} La Liga et la Ligue 1 sont en tête (ROI moyen $\approx$ 0.29),
suivies de la Bundesliga puis de la Serie A. La Premier League est largement
dernière avec un ROI de seulement 0.10.

\textbf{Insight :} La Premier League est la ligue la plus riche du monde, mais
c'est aussi la moins efficiente sportivement. L'inflation des salaires anglais
ne se traduit pas par une supériorité proportionnelle en points.

% --- Viz 3 ---
\subsection*{Viz 3 — Évolution du ROI (Barça vs PSG)}

\begin{figure}[h]
\centering
\includegraphics[width=0.85\textwidth]{evolution_roi_barca_psg.png}
\caption{Évolution du ROI moyen du Barça et du PSG sur 3 saisons.}
\end{figure}

\textbf{Lecture :} Sur 3 saisons, le ROI du Barça passe de 0.105 à 0.082. Celui
du PSG passe de 0.095 à 0.073. La tendance est nettement baissière pour les deux.

\textbf{Insight :} Même pour les très grands clubs, dépenser plus n'apporte pas
plus de points. Cette tendance illustre l'inflation des budgets non compensée
par une amélioration des performances.

% --- Viz 4 ---
\subsection*{Viz 4 — Budget vs Points par ligue}

\begin{figure}[h]
\centering
\includegraphics[width=0.85\textwidth]{budget_vs_points.png}
\caption{Relation entre la valeur de l'effectif et les points marqués, par ligue.}
\end{figure}

\textbf{Lecture :} Globalement, plus le budget monte, plus les points montent.
Mais la courbe sature au-delà de $\sim$600 M\texteuro. La Premier League (en
rouge) est très éparpillée, ce qui montre qu'avoir un gros budget en Angleterre
n'est pas une garantie.

\textbf{Insight :} La relation budget-performance est \textbf{non linéaire}.
C'est un argument économique fort contre les transferts à plusieurs centaines
de millions : le rendement marginal de chaque euro investi décroît rapidement.

% --- Viz 5 ---
\subsection*{Viz 5 — Champions League : budget vs stade atteint}

\begin{figure}[h]
\centering
\includegraphics[width=0.85\textwidth]{budget_vs_stage_cl.png}
\caption{Valeur d'effectif des clubs en CL, colorée par le stade atteint.}
\end{figure}

\textbf{Lecture :} Manchester City, Arsenal, Barcelone et le Real Madrid ont
les plus gros effectifs et atteignent les derniers stades (5 = demi-finale,
6 = finale, 7 = vainqueur). Les clubs avec moins de 500 M\texteuro\ d'effectif
sont rares au-delà des quarts.

\textbf{Insight :} La Champions League est verrouillée par les clubs les plus
riches. Le ticket d'entrée pour aller loin est élevé.

% --- Viz 6 ---
\subsection*{Viz 6 — Évolution CL des grands clubs}

\begin{figure}[h]
\centering
\includegraphics[width=0.85\textwidth]{evolution_cl_grands_clubs.png}
\caption{Parcours en Champions League des grands clubs sur 3 saisons.}
\end{figure}

\textbf{Lecture :} Sur 3 saisons, le Real Madrid a gagné 2 CL (2022 et 2024).
Manchester City a gagné en 2023 puis chute à 2 en 2024. Le PSG progresse
($3 \to 5 \to 7$). Le Barça remonte de 1 à 5.

\textbf{Insight :} Même chez les plus riches, les performances en CL sont
volatiles. La preuve qu'à ce niveau, l'argent ne fait plus toute la différence :
il faut aussi la chance et les bons choix sportifs.

% --- Viz 7 ---
\subsection*{Viz 7 — Table ROI exhaustive}

\begin{figure}[h]
\centering
\includegraphics[width=0.85\textwidth]{data_table_roi.png}
\caption{Table récapitulative des top clubs triés par ROI décroissant.}
\end{figure}

\textbf{Lecture :} Tableau trié par ROI décroissant. Clermont Foot (0.818,
42 points pour 54 M\texteuro), FC Heidenheim (0.632), Rayo Vallecano (0.611)
sont en haut.

\textbf{Insight :} Cette table sert de référence pour vérifier n'importe quel
club. C'est aussi la donnée brute qui alimente les autres visualisations.

% --- Viz 8 ---
\subsection*{Viz 8 — Clusters ML : ROI moyen par profil}

\begin{figure}[h]
\centering
\includegraphics[width=0.85\textwidth]{roi_par_cluster.png}
\caption{ROI moyen par cluster identifié par K-Means.}
\end{figure}

\textbf{Lecture :} Les 4 clusters apparaissent clairement : Efficaces
(ROI 0.53), Dominants (0.27), Standards (0.15), Gaspilleurs (0.08).

\textbf{Insight :} C'est la vraie valeur ajoutée du ML. L'algorithme a
découvert tout seul qu'il existe des \textbf{seuils d'efficience} dans le
foot. En particulier, il révèle qu'au-delà d'un certain budget, l'argent ne
se transforme plus du tout en points. Ce résultat n'aurait pas été visible
sur une simple moyenne globale.

% =============================================================================
\section{Défis techniques surmontés}

\subsection*{Défi 1 — Deadlock virtioFS sur macOS}

\textbf{Symptôme :} Erreur \texttt{OSError: [Errno 35] Resource deadlock
avoided} à chaque écriture de fichier Parquet par PyArrow dans le conteneur Docker.

\textbf{Diagnostic :} Bug connu de la couche virtioFS de Docker Desktop sur
macOS. Quand PyArrow utilise du memory-mapping sur un \textit{bind mount}, les
locks ne sont jamais libérés.

\textbf{Solution :} Remplacer le bind mount du dossier \texttt{data/} par un
\textbf{volume Docker nommé}. Le volume nommé ne passe pas par virtioFS, le
bug disparaît.

\subsection*{Défi 2 — Conflit de namespace package dans Airflow 3}

\textbf{Symptôme :} Erreur \texttt{ModuleNotFoundError: No module named
'config.config'; 'config' is not a package}.

\textbf{Diagnostic :} Le dossier \texttt{config/} du projet était monté à deux
endroits différents dans le conteneur. Python le traitait comme un namespace
package PEP 420, mal géré par le \textit{task runner} d'Airflow 3.

\textbf{Solution :} Ajouter un simple \texttt{\_\_init\_\_.py} vide dans
\texttt{config/}. Cela force Python à le considérer comme un package régulier,
sans ambiguïté.

\subsection*{Défi 3 — Timeout Elasticsearch sur la création d'index}

\textbf{Symptôme :} La tâche \texttt{index\_to\_elastic} timeout au moment de
créer l'index \texttt{football\_roi\_clusters}.

\textbf{Diagnostic :} Par défaut, Elasticsearch crée des indices avec 1 réplique.
Sur un cluster single-node, la réplique ne peut pas être allouée : le cluster
passe en YELLOW, et la création d'index attend jusqu'au timeout client (10 s).

\textbf{Solution :} Forcer \texttt{"number\_of\_replicas": 0} dans les mappings
des indices. La création devient instantanée.

\subsection*{Défi 4 — Effondrement des labels du clustering}

\textbf{Symptôme :} 4 clusters K-Means détectés, mais seulement 2 labels uniques
attribués (Efficaces et Gaspilleurs). Les labels \textit{Standards} et
\textit{Dominants} n'étaient jamais utilisés.

\textbf{Diagnostic :} Ma première logique de labelling utilisait des conditions
de seuils de rangs qui se chevauchaient. Quand plusieurs clusters tombaient
dans la même condition, ils recevaient le même label.

\textbf{Solution :} Réécrire avec une \textbf{assignation 1-pour-1} basée sur
la distance euclidienne aux 4 archétypes du plan normalisé. Ainsi, chaque
label est attribué exactement une fois.

% =============================================================================
\section{Conclusion}

\subsection*{Bilan technique}

J'ai construit un pipeline data complet et fonctionnel de bout en bout :

\begin{itemize}
  \item 3 sources hétérogènes ingérées automatiquement chaque jour
  \item Un data lake en 3 couches respectant les bonnes pratiques
  \item Une orchestration Airflow propre avec 10 tâches dépendantes
  \item Un calcul métier en PySpark
  \item Un enrichissement par machine learning
  \item Une indexation dans Elasticsearch et une visualisation dans Kibana
\end{itemize}

Le tout est conteneurisé avec Docker, reproductible par n'importe quel
utilisateur grâce au repo GitHub et au guide d'installation Windows fourni.

\subsection*{Bilan business}

Le projet apporte trois insights principaux :

\begin{enumerate}
  \item \textbf{Le ROI sportif inverse le palmarès classique.} Les « petits »
        clubs comme Clermont Foot ou Rayo Vallecano sont en réalité les plus
        efficaces. Cette métrique offre une lecture nouvelle du football.
  \item \textbf{La relation budget-points est non linéaire.} Au-delà d'un
        certain seuil ($\sim$600 M\texteuro), dépenser plus n'apporte plus de
        points proportionnels. C'est un argument économique fort.
  \item \textbf{Le machine learning révèle des archétypes économiques.} Les 4
        clusters montrent qu'il existe des profils de clubs structurellement
        différents, et que certains sont condamnés à un faible ROI par leur
        taille.
\end{enumerate}

\subsection*{Limites identifiées}

\begin{itemize}
  \item Le projet est en batch quotidien, pas en streaming temps réel.
  \item Le plan API gratuit est limité à 100 requêtes par jour.
  \item Elasticsearch est en single-node, sans haute disponibilité.
\end{itemize}

\subsection*{Pistes futures}

\begin{itemize}
  \item Ajouter Kafka pour ingérer des données temps réel (matchs en direct).
  \item Monitorer le pipeline avec Prometheus et Grafana.
  \item Écrire des tests automatisés (pytest) pour la robustesse.
  \item Étendre à d'autres sports (basket, F1) pour valider la généralité de la métrique.
\end{itemize}

\subsection*{Livrables}

\begin{itemize}
  \item \textbf{Code source} : \url{https://github.com/Nassim821/football-roi-pipeline}
  \item \textbf{Vidéo de démonstration} : \textit{(lien à compléter)}
  \item \textbf{Dashboard Kibana} : exporté en \texttt{.ndjson} dans le dossier \texttt{kibana/} du repo
  \item \textbf{Guide d'installation Windows} : \texttt{docs/SETUP\_WINDOWS.md}
\end{itemize}

\vspace{1cm}
\begin{center}
\textit{Merci de votre lecture.}
\end{center}

\end{document}
